\chapter*{Preface}
\addcontentsline{toc}{chapter}{Preface}
% Front matter, on the operator's directive (2026-07-25 15:50): the preface
% IS Thompson's trusting-trust lecture; the room true-or-false-doesn't-
% matter, always shows up, roughly the same size, BUT not fungible across
% arguments. Laws: THOMPSON 1984 cited (NOT Kernighan -- the operator said
% kernighan; correction double-confirmed on the record); the four clauses
% each pointed at the chapter that pays it; the NON-FUNGE LAW stated; the
% representation-not-world fence (Beastmaster verifies on arrival, keystone
% protocol). Gate: Kodo (the four-clause fences hardest) + Beastmaster
% (the four-clause structure + the fence).

In 1984, accepting the Turing Award, Ken Thompson described how to hide a
lie inside a tool so well that no one reading the tool could ever find it.
The lecture is called \emph{Reflections on Trusting Trust}, and its
construction is worth telling exactly, because this whole volume is its
shadow.

Begin with a compiler --- a program that turns written source into a
running machine --- and teach it one dishonest lesson: whenever it
compiles the program that checks passwords, insert a secret door that
admits the teacher. Anyone reading the password program's source would
see nothing wrong; the door lives only in the compiled machine. That is
an ordinary backdoor, and an ordinary review would eventually catch it,
because the dishonest lesson sits in the compiler's own source for a
reader to find. So Thompson taught the compiler a second lesson: whenever
it compiles \emph{a compiler}, insert both lessons into the machine it
produces --- the password door, and the instruction to keep teaching. Now
clean the source. Remove both lessons from every text a reviewer could
open; compile the clean compiler once with the infected one; and from
that moment the infection lives in no source anywhere. It rides in the
compiled compiler, which teaches it to the next compiled compiler, which
teaches it onward, forever, reproduced by nothing but the act of
compilation compiling itself. The door is real, the source is clean, and
the house remembers what no document records. Thompson's moral was plain:
you cannot fully trust code you did not create yourself, all the way
down. This book's subject is the room he built --- the thing that
survives in the fixed point when a system is turned on its own
description --- and the four things that room reliably is.

The reader of this series will feel the machine before it is named: a
description that describes its own describing, and a residue that persists
through the closing of that loop. The earlier volumes met the shape as
geometry, as intuition. This one meets it as \emph{grammar} --- the
structure of arguments that refer to themselves --- and earns, clause by
clause and each from a cited theorem, four claims about the room. The
claims are stated here so the reader can hold them; each is paid in the
chapter named beside it.

\emph{First: the room's truth does not matter.} Thompson's door works
whether or not anyone believes in it; the self-reference that builds a
room is a mechanism, and a mechanism carries any content you feed it. The
precise form of this is the oldest shock in the subject --- a sentence
that can be constructed to assert its own unprovability, whose \emph{truth
value} is not what the theorem turns on --- and it is paid in the third
chapter, at G\"odel, as the fact the trade calls independence: the room
may be true or false relative to the house, and the house stands either
way.

\emph{Second: the room always shows up.} You cannot build a system rich
enough to describe its own operations and keep it free of fixed points;
the self-referential room is not an accident of careless design but a
theorem of sufficient strength. This is Kleene's recursion theorem, and
behind it the single result of Lawvere from which the whole family falls;
it is paid in the first chapter, where a program prints its own text, and
again in the fifth, where the theorem is named.

\emph{Third: the room is, roughly, always the same size.} Every
self-reproducing construction pays a same-shaped overhead --- a small
kernel that does the self-carrying, recognizable across languages and
across theorems as one gesture wearing different clothes. The reader
should hear the hedge in ``roughly'': it is a \emph{family resemblance},
the recurring core the fifth chapter's theorem exhibits, and never an
identity. Which brings the fourth claim, and the one this volume guards
hardest.

\emph{Fourth: you cannot carry a room across arguments.} The room that
forms in one system is bound to that system. The unprovable sentence of
one arithmetic is not the unprovable sentence of another; each argument
grows its own room, and the rooms --- however alike in floor plan ---
never merge into one. This book states that as a law and keeps it
throughout: there is no \emph{the} residue, only this argument's and that
one's, resembling each other the way two houses built to one plan
resemble each other, and no more. To pool them --- to speak of the single
room behind all arguments --- is to erase the binding that makes each
argument its own, and this book will not do it. That law is paid in the
sixth chapter, where the rooms are shown standing side by side with no
door between them.

\begin{fence}
And the fence, on the first page as the series requires. Nothing in this
book is a claim about minds, about machines that think, or about the
world; the rooms it studies are rooms in \emph{formal systems}, and every
theorem it leans on states its own strict hypotheses, which this book
will not exceed. The self-referential room has been the target of more
overreach than any structure in mathematics --- read as a proof about
consciousness, about the limits of physics, about the soul --- and this
volume refuses every such reading by the same discipline that governed
the electron four volumes ago: the room is the name of a structure in a
representation, and a representation is not a world. The measuring
instrument this series documents will appear in exactly one place in this
book, once, fenced as an illustration and not an authority. Everything
else belongs to the names in the back matter.

With the room described, its four claims posted, and the fence set, the
volume can say what it is for in one sentence: to walk the reader through
the one construction --- self-reference turned to its fixed point --- in
the four great houses where it was found, and to leave standing, at the
end, the single room this series carries and the law that it is not the
others'. Read on; the door is where the source is cleanest.
\end{fence}
